\documentclass[11pt,a4paper]{article}
\usepackage{fontspec}
\setmainfont{TeX Gyre Termes}[Ligatures=TeX]
\usepackage{amsmath,amssymb,bm}
\usepackage{booktabs}
\usepackage{graphicx}
\usepackage{geometry}
\geometry{margin=2.3cm}
\usepackage[hidelinks]{hyperref}
\usepackage{caption}
\captionsetup{font=small}
\usepackage{enumitem}
\usepackage{mdframed}

\newcommand{\cL}{\mathcal{L}}
\newcommand{\cM}{\mathcal{M}}
\newcommand{\cN}{\mathcal{N}}

\title{\textbf{A Two-Stage Procedure for Consolidating and Re-Mapping an AI Risk Taxonomy}\\[3pt]
\large Stage 1: level-constrained L4 consolidation at three resolutions.\\
Stage 2: damped, seed-anchored EM assignment to a fixed L3 layer}
\author{RAI Risk Taxonomy 2.0 --- technical report (rev.~3)}
\date{August 10, 2026}

\begin{document}
\maketitle

\begin{abstract}
We report a completed two-stage procedure that (i) reduces an inventory of 1{,}652 active L4 risk cards to three nested release resolutions and (ii) re-maps every resulting card onto a fixed 56-node L3 layer. Stage 1 combines a single removal criterion (absence of harm content), a form-only normalization that freezes each card's harm mechanism and outcome, a pre-approved exact-duplicate merge dictionary, and constrained complete-linkage consolidation cut at $\tau=0.80$ and $\tau=0.70$; a nameability gate and a cross-abstraction-level merge prohibition prevent the umbrella-card failure observed in an earlier round. The resulting sets are Master 1{,}612, C80 1{,}369, and C70 937 cards, with zero cannot-link violations, zero cross-level merges, exact nesting, and lossless ID mapping. Stage 2 replaces unconstrained EM --- which we show is \emph{not identified} on this data (independent random initializations agree on only 0.3--3.3\% of assignments) --- with a damped, seed-anchored variant that keeps class centroids within cosine 0.85 of enriched L3 seeds. After a two-stage assignment (automatic, then forced with expert adjudication of 814 low-margin holds), all three sets use 54 of 56 L3 nodes, both taxonomy-hold nodes are emptied, and sensitivity across three axes with two repetitions each yields 92--99\% agreement.
\end{abstract}

\section{Problem and prior failure}

Two structural defects motivated the work. First, the L4 inventory had accumulated near-duplicate cards from multiple source repositories, with heterogeneous writing conventions: roughly half of all definitions were descriptive prose or verbatim excerpts rather than risk definitions. Second, 594 cards (36.8\%) sat in \emph{taxonomy-hold} L3 nodes, i.e., were effectively unclassified, while several substantive L3 nodes held one or two cards.

An earlier consolidation round failed in an instructive way. Merging cards without regard to abstraction level produced umbrella cards: for example, six instance-level cards were absorbed into a construct-level card and renamed ``over-reliance and automation bias.'' Diagnostically, construct-level cards behave as \emph{merge magnets} --- they are uniformly similar to an entire family, so they pull heterogeneous members into one cluster. We measured this: the number of neighbours above $\cos \geq 0.80$ has median 0 across the inventory but reaches 8--14 for construct cards such as ``diminished human oversight'' and ``lack of transparency and explainability.'' The corrective is not to rewrite such cards --- established risk names (automation bias, goal misalignment, loss of control) must survive verbatim for citation and retrieval --- but to forbid merges across levels.

\section{Stage 1: consolidation}

\subsection{Setup}

Let $\cM_0$ be the 1{,}652 active cards. Card $r_i$ is embedded from
\begin{equation}
x_i = \texttt{label\_en}_i \texttt{.}\; \texttt{definition\_en}_i \;\texttt{/}\; \texttt{label\_ko}_i \texttt{.}\; \texttt{definition\_ko}_i ,
\qquad
\mathbf{e}_i = \frac{\mathrm{Enc}(x_i)}{\lVert\mathrm{Enc}(x_i)\rVert_2} \in \mathbb{S}^{1023},
\end{equation}
with $\mathrm{Enc}$ = BGE-M3 (\texttt{max\_seq\_length}=256). Similarity is $s_{ij}=\mathbf{e}_i^\top\mathbf{e}_j$, distance $d_{ij}=1-s_{ij}$.

Two protective structures are defined before any operation. The \textbf{cannot-link set} $\cN$ holds seven expert-adjudicated pairs that must never merge, including embedding false positives ($\cos$ up to 0.944) and pairs whose conflation would dilute a distinct harm (child sexual exploitation vs.\ adult sexual content; moral responsibility gap vs.\ legal liability gap). The \textbf{protected-term list} (44 entries drawn from TTA, KISA, NIS, ETRI and the Korean AI Framework Act, minus generic nouns) marks labels that may not be renamed at any step.

Each card is tagged \emph{construct} or \emph{instance}: instance iff its definition binds a context (application domain, affected party, decision type, or system type). This is metadata only; no card is rewritten to change its level. The inventory splits 1{,}218 construct / 434 instance.

\subsection{Filtering and normalization}

Removal uses exactly one criterion: the definition contains no harm content (literature-classification memos, section preambles, pure concept definitions, pure countermeasures). Abstractness, missing context, and capability phrasing are \emph{not} removal grounds --- dangerous-capability items (self-replication, deception, offensive cyber capability) are standard risk entries in the AI-safety literature and are retained. Three review rounds over the union of prior candidates produced \textbf{40 removals}; 25 earlier candidates were restored under the narrowed criterion.

Normalization is form-only under three invariants: the harm mechanism and outcome are frozen, no context may be added (a construct card stays construct), and protected terms may not be renamed. 1{,}335 cards were normalized (786 label + definition, 549 definition only) and audited by two senior reviewers, who caught four mechanism substitutions --- e.g.\ a rewrite that turned ``the model generates harmful content'' into ``the safety filter fails,'' replacing a generation mechanism with a control-failure mechanism --- and three cases of context injected into construct cards.

\subsection{Consolidation}

Consolidation is agglomerative: repeatedly merge the most similar pair of live units and recompute. The only degree of freedom is the inter-unit similarity $\sigma$. We use complete linkage, $\sigma_{\mathrm{comp}}(\cL_a,\cL_b)=\min_{i\in\cL_a,j\in\cL_b} s_{ij}$, on three empirical grounds: at $\tau=0.70$ centroid and single linkage collapse (largest cluster 1{,}191 and 1{,}514 cards, versus 8 for complete); only complete linkage propagates the $d_{ij}=\infty$ encoding of cannot-link constraints; and complete linkage guarantees every within-cluster pair exceeds $\tau$, which is what a merged card claims.

Two hard constraints enter the distance matrix: $d_{ij}\leftarrow\infty$ for $(i,j)\in\cN$, and $d_{ij}\leftarrow\infty$ whenever $r_i$ and $r_j$ have different abstraction levels. The second constraint is the corrective for the umbrella failure.

A \emph{merge dictionary} is applied before clustering. Exact textual duplicates and adjudicated high-similarity pairs are merged with three decisions taken independently: the surviving ID, the final label, and the final definition --- chosen by card quality, not by which card the similarity search treated as the neighbour. Chain formation is blocked: taking the transitive closure of pairwise merges produced a 22-card chain in the loss-of-control family, so groups are assigned exclusively in descending confidence and conflicting groups are deferred to the $\tau$ cut, where complete linkage cannot chain. The dictionary contributes 120 groups (258 $\to$ 120 cards).

One dendrogram is built and cut at $\tau=0.80$ and $\tau=0.70$. By monotonicity of complete linkage this yields $\cM \supseteq \mathrm{C80} \supseteq \mathrm{C70}$, so cross-tier ID mapping is a tree.

\subsection{Nameability gate}

Prior threshold analysis on a 12-pair golden set placed the safe band at $0.84 \leq \tau < 0.94$; both target resolutions lie below it. We therefore do not promote $\tau$-clusters directly. For each candidate cluster two reviewers independently attempt a genuine hypernym as a noun phrase of at most twelve English words. Approval requires success and agreement; rejection follows from enumerative names, empty umbrella terms, or heterogeneous members. A third criterion enforces the level constraint at the naming step: \emph{if the only adequate name is an established construct name, the merge is rejected}, since needing that name is evidence the cluster has been flattened upward. This rule fired repeatedly --- clusters that would have collapsed into ``loss of control,'' ``goal misalignment,'' ``responsibility gap,'' or ``explainability'' were rejected. Rejections propagate from C70 to C80 by member set, preserving nesting without re-clustering. Outcome: C70 325/451 approved, C80 127/157 approved.

\subsection{Stage 1 result}

\begin{table}[h]
\centering\small
\begin{tabular}{@{}lrrr@{}}
\toprule
Set & Cards & vs.\ Master & Merges \\
\midrule
Source inventory & 1{,}652 & --- & --- \\
Master (after 40 removals, 1{,}335 normalizations) & 1{,}612 & --- & --- \\
C80 ($\tau=0.80$) & 1{,}369 & $-15.1\%$ & 172 \\
C70 ($\tau=0.70$) & 937 & $-41.9\%$ & 366 \\
\bottomrule
\end{tabular}
\end{table}

Audit of all three sets: lossless ID mapping (JSON and CSV), zero cannot-link violations, \textbf{zero cross-level merges}, zero nesting violations, zero duplicate labels in either language, zero label-length violations, zero definition-format violations. Largest cluster 7 (C80) and 12 (C70), against 22 in the failed round.

\section{Stage 2: mapping to the fixed L3 layer}

\subsection{Unconstrained EM is not identified}

The release pipeline had historically used EM over L3 seed centroids. Running it under a sensitivity design (three axes, two repetitions) exposed a fundamental problem: independent random initializations converge to comparable objective values but agree on only \textbf{0.3--3.3\%} of assignments. The partition is determined by the initialization, not by the data. The cause is the M-step: unconstrained centroid updates drift away from the L3 seeds, with minimum centroid--seed cosine falling to 0.605 and agreement with pure nearest-seed assignment at 46\%. The nodes still carry L3 labels but their centroids no longer represent those L3s.

\subsection{Damped, seed-anchored EM}

We replace the M-step with a damped update,
\begin{equation}
\bm{\mu}_k \;\leftarrow\; \frac{\alpha\,\mathbf{s}_k + (1-\alpha)\,\overline{\mathbf{e}}_k}{\lVert \alpha\,\mathbf{s}_k + (1-\alpha)\,\overline{\mathbf{e}}_k \rVert}, \qquad \alpha = 0.4,
\end{equation}
where $\mathbf{s}_k$ is the L3 seed and $\overline{\mathbf{e}}_k$ the current member mean; empty classes revert to the seed. Initialization is always the seed --- random initialization is recorded in the manifest as an \emph{inadmissible specification}, not a sensitivity condition. Seeds are enriched: each non-hold L3 seed text is augmented with the labels of its eight nearest previously-assigned cards, strengthening the anchor. Hold nodes are masked out of the assignment space entirely, so no card can be assigned to a taxonomy-hold node.

The $\alpha$ sweep is reported in the notebook: $\alpha=0$ gives 46\% agreement with nearest-seed and centroid drift 0.605; $\alpha=0.4$ gives 0.848; $\alpha=0.7$ gives 0.965 at the cost of a 302-card largest class. $\alpha=0.4$ balances anchoring against class balance.

\subsection{Two-stage assignment}

Stage 2a assigns automatically and holds any card with margin $m_i = s_{(1)}-s_{(2)} < 0.03$. Stage 2b resolves holds by expert adjudication and forces an assignment for every remaining card, so no card ends unassigned. For Master, 814 holds were adjudicated card-by-card against the full 56-node catalogue (not merely the top-3 candidates), with hold nodes prohibited; confidence was high 188, medium 506, low 120. The tiers inherit from Master --- a tier card takes the L3 of its Master sources --- which guarantees cross-tier hierarchy consistency; only the 191 ties (C80 53, C70 138) were adjudicated separately, judged by the merged card's own harm mechanism rather than by member votes.

\subsection{Sensitivity}

Three axes, two repetitions each (repetition is a 90\% bootstrap subsample, since damped EM is deterministic given the seed initialization): card text (bilingual vs.\ English-only), L3 seed text (enriched vs.\ plain), and damping ($\alpha=0.4$ vs.\ $0.6$).

\begin{table}[h]
\centering\small
\begin{tabular}{@{}lrrrr@{}}
\toprule
Condition & \multicolumn{2}{c}{Stage 2a (agreement, ARI)} & \multicolumn{2}{c}{Stage 2b} \\
\cmidrule(lr){2-3}\cmidrule(lr){4-5}
 & all & --- & all & free cards \\
\midrule
Axis 1: English-only text & 75.2\% & 0.57 & 96.1\% & 92.1\% \\
Axis 2: plain L3 seed     & 69.9\% & 0.50 & 96.1\% & 92.1\% \\
Axis 3: $\alpha=0.6$      & 89.7\% & 0.79 & 98.6\% & 97.3\% \\
Bootstrap 90\% (2 seeds)  & 91--99\% & --- & 99.4/99.7\% & 98.9/99.3\% \\
\midrule
Unconstrained EM, random init (rejected) & 0.3--3.3\% & 0.19--0.29 & --- & --- \\
\bottomrule
\end{tabular}
\caption{Stage 2b figures fix the 811 adjudicated cards and let 801 float; ``free cards'' reports agreement on the floating subset only. The notebook's raw Stage 2b output shows 100\% because the decision file pinned every card, leaving EM no role; those numbers are void and the figures above supersede them.}
\end{table}

Minimum centroid--seed cosine rises from 0.605 (unconstrained) to 0.848--0.915 (damped), confirming that the classes remain the L3s they are labelled as.

\subsection{Stage 2 result}

All three sets assign every card, use \textbf{54 of 56} L3 nodes, and leave both taxonomy-hold nodes empty. The 594 previously unclassified cards are distributed across substantive L3s; 1{,}108 of 1{,}612 Master cards changed assignment, which is expected since 1{,}335 cards had their text normalized and the prior assignment was computed on the pre-normalization text. No substantive L3 is empty. Largest class 90 (governance-framework absence), smallest 6 (hardware and mechanical faults). Single-source tier cards agree with their Master assignment in every case.

\section{Limitations}

The nameability gate is a procedural, not statistical, guarantee, and the golden set behind the $[0.84,0.94)$ band contains only twelve pairs, each boundary set by a single observation. Sensitivity is confined to a single embedding model: alternate models could not be retrieved through the corporate TLS-intercepting proxy, so model sensitivity --- the strongest available test --- is untested. Seed enrichment draws exemplars from prior human assignments, which introduces a mild circularity that we treat as curated-exemplar seeding and expose as sensitivity axis 2 (agreement with plain seeds is 92.1\% on free cards). Severity aggregates by maximum and probability by maximum, the latter a lower bound for a union of events and therefore anti-conservative for risk management; member-level values are preserved so the aggregate can be revised. Finally, C70 retains one single-card L3 under the fixed-L3 policy.

\section{Reproducibility}

Both stages run as notebooks that hash every input, record package versions, cache embeddings under a text hash so that any edit invalidates the cache, and fingerprint the clustering so that stale human-decision files are rejected rather than silently applied. Artifacts: \texttt{data/experiments/stage1/out/} (master, c80, c70, mappings, audit) and \texttt{data/experiments/stage2/out/} (assignments, mappings, sensitivity tables, manifest).

\begin{thebibliography}{9}
\small
\bibitem{slattery2024} P.~Slattery et al. The AI Risk Repository: A comprehensive meta-review, database, and taxonomy of risks from artificial intelligence. arXiv:2408.12622, 2024.
\bibitem{weidinger2022} L.~Weidinger et al. Taxonomy of risks posed by language models. In \emph{ACM FAccT}, 2022.
\bibitem{shelby2023} R.~Shelby et al. Sociotechnical harms of algorithmic systems: Scoping a taxonomy for harm reduction. In \emph{AAAI/ACM AIES}, 2023.
\bibitem{nist2023} NIST. Artificial Intelligence Risk Management Framework (AI RMF 1.0). NIST AI 100-1, 2023.
\bibitem{chen2024bgem3} J.~Chen et al. BGE M3-Embedding: Multi-lingual, multi-functionality, multi-granularity text embeddings through self-knowledge distillation. arXiv:2402.03216, 2024.
\bibitem{abbas2023} A.~Abbas et al. SemDeDup: Data-efficient learning at web-scale through semantic deduplication. arXiv:2303.09540, 2023.
\bibitem{hastie2009} T.~Hastie, R.~Tibshirani, J.~Friedman. \emph{The Elements of Statistical Learning}, 2nd ed. Springer, 2009.
\bibitem{wagstaff2001} K.~Wagstaff, C.~Cardie, S.~Rogers, S.~Schr\"odl. Constrained k-means clustering with background knowledge. In \emph{ICML}, 2001.
\bibitem{traag2019} V.~A.~Traag, L.~Waltman, N.~J.~van Eck. From Louvain to Leiden: guaranteeing well-connected communities. \emph{Scientific Reports}, 9:5233, 2019.
\bibitem{shah2022} R.~Shah et al. Goal misgeneralization: Why correct specifications aren't enough for correct goals. arXiv:2210.01790, 2022.
\bibitem{carlsmith2022} J.~Carlsmith. Is power-seeking AI an existential risk? arXiv:2206.13353, 2022.
\bibitem{matthias2004} A.~Matthias. The responsibility gap: Ascribing responsibility for the actions of learning automata. \emph{Ethics and Information Technology}, 6(3):175--183, 2004.
\bibitem{biggio2013} B.~Biggio et al. Evasion attacks against machine learning at test time. In \emph{ECML PKDD}, 2013.
\bibitem{tta} TTA. Information and Communication Terminology Dictionary (adversarial AI, deepfake, explainable AI). Telecommunications Technology Association, Korea.
\bibitem{kisa2026} MSIT and KISA. AI Security Red Teaming Guide, 2026.
\bibitem{etri2024} ETRI. Standardization trends in advanced AI safety and trustworthiness. \emph{Electronics and Telecommunications Trends}, 39(5), 2024.
\bibitem{aiact2026} Republic of Korea. Framework Act on the Development of Artificial Intelligence and Establishment of a Foundation of Trust, Act No.~20676, in force 22 January 2026.
\end{thebibliography}

\end{document}
